
\subsubsection{6.1 Interpretation of the Advantage Variance Finding}

The 76× advantage variance gap between function-level and repository-level GRPO training is not attributable to hyperparameter differences, training duration differences, or model initialization differences—all of these are held constant. The gap is explained algebraically: when positive reward rate ≈0\%, std(r)=0 for essentially every training group, and the GRPO gradient is zero by definition. The result confirms that binary execution feedback on competitive programming problems is insufficient to produce informative GRPO gradients for CodeLlama-7b-Instruct-hf.

This finding does not imply that GRPO is generally ineffective. The repository-level condition (SWE-bench Verified) maintains advantage variance of 0.317, which indicates non-degenerate training. The distinction is between the reward structure (binary vs. partial credit) and the model's capability level relative to the task difficulty, not the algorithm itself.

\subsubsection{6.2 Implications for Curriculum GRPO}

The curriculum GRPO hypothesis posits that starting training on easier problems should maintain higher reward density during early training, preserving informative advantage groups and enabling the policy to improve before encountering harder problems. The H-M1 finding provides mechanistic motivation for this hypothesis: if early-curriculum problems yield non-zero reward rates, advantage groups are non-degenerate, and learning can occur.

However, the H-E1 smoke test showed reward density = 0.0 for all four conditions including easy_only. This raises a question about the activation condition for curriculum GRPO: if even easy-tier APPS problems yield zero reward for the model at initialization, easy-to-hard curriculum ordering provides no gradient advantage over uniform sampling, because both yield degenerate advantage groups throughout early training. Confirming or ruling out this possibility requires measuring per-difficulty-tier solve rates at initialization before committing to full training runs.

\subsubsection{6.3 Limitations}

\textbf{L1: Cross-granularity, not within-difficulty comparison.} H-M1 compares function-level (APPS+CodeContests) vs. repository-level (SWE-bench Verified) advantage variance. This establishes that reward sparsity causes advantage collapse but does not test whether easy-tier APPS problems yield higher reward density than hard-tier APPS problems within the function-level regime. The curriculum hypothesis requires a within-difficulty-level comparison that was not executed.

\textbf{L2: Single model.} The advantage variance comparison used CodeLlama-7b-Instruct-hf. Results may differ for other 7B architectures, for base models vs. instruction-tuned models, or for models at different capability levels. The finding is likely to generalize to other 7B-class models on competitive programming, but this has not been verified.

\textbf{L3: Behavioral hypothesis untested.} The primary behavioral claim—that easy-to-hard curriculum ordering improves HumanEval+ pass@1 relative to uniform sampling—was not tested. No full 5000-step training run was executed, and no pass@1 evaluation data exists for any condition. The paper's empirical contribution is the mechanistic characterization of advantage collapse (H-M1) and the infrastructure validation (H-E1), not a demonstration of curriculum effectiveness.

\textbf{L4: Assumption A1 unconfirmed.} The curriculum mechanism depends on the initializing policy having a non-zero solve rate on easy-tier problems (APPS tiers 0–2). The H-E1 smoke test showed 0.0 reward density even for the easy_only condition at initialization. If assumption A1 fails, the curriculum mechanism has no activation pathway. A targeted evaluation measuring per-tier solve rates at initialization is needed before proceeding to full training.

\textbf{L5: 120-step measurement window.} The advantage variance measurement was conducted over 120 training steps. This window is sufficient to characterize the reward distribution at initialization but does not capture whether variance might change over longer training horizons as the policy adapts.

\textbf{L6: Reward function confound.} The two conditions differ in both task granularity and reward function structure (binary vs. partial credit). It is not possible from this experiment alone to attribute the variance difference solely to task granularity rather than reward function design. Partial credit rewards applied to function-level tasks could potentially produce non-zero advantage variance even on competitive programming problems.

\subsubsection{6.4 Practical Implications}

Per-step advantage variance, measurable over as few as 10–20 training steps, provides a diagnostic signal for whether GRPO will produce informative gradients under a given reward function and model capability combination. A near-zero average advantage variance in early steps indicates that the reward structure is too sparse for the current model, regardless of other training configuration choices. This diagnostic is low-cost compared to full training runs.
